\begin{center}
\textbf{{\Large ABSTRACT}}\\
\begin{center}
\justifying
Linear recurrent neural networks offer an efficient alternative to attention for long sequences, but their ability to track state depends on the structure of their state-transition matrix. DeltaProduct constructs this matrix at every time step as a product of $n_h$ generalized Householder factors, where $n_h$ is a single hyperparameter chosen by hand and applied identically to every token. Increasing $n_h$ increases both the expressivity and the computation spent on every token, whether or not that token requires it. This project studies how many factors each token actually requires, and replaces the fixed global order with a per-token order that is learned.

The base method is first reproduced on the $S_3$ and $S_5$ word problems on a single 6~GB GPU, reaching a token accuracy of 0.864 at sequence length 512 after training at length 128. We then show, by explicit construction verified to machine precision, that the minimal number of factors a token requires is not its transposition length but the rank of its image under the cheapest faithful representation of the group generated by the entire input alphabet. A 5-cycle therefore requires four factors when the alphabet generates $S_5$, but only two when it generates $A_5$, whose three-dimensional rotation representation is unavailable to $S_5$. Using a reference allocation derived from this analysis, per-token order reaches the accuracy of the maximal fixed order while requiring between 1.15 and 1.75 factors per token instead of four.

We design a monotone halting gate that scales the strength of each factor per token and train it end to end with a scheduled compute-budget penalty and no supervision on the allocation. In three of ten runs the learned gate recovers the demand structure exactly, assigning two factors to low-demand tokens and four to high-demand tokens with no under-allocation, at 2.20 required factors per token and a token accuracy of 0.9998, compared with 0.9696 for a fixed three-factor model. Training on these tasks is bistable across runs, so results are reported as the fraction of repeated runs that succeed. The reported savings are measured in required factors per token; converting them into wall-clock savings requires ragged execution of the factor product, which is identified as the principal direction for future work.
\end{center}
\end{center}
\noindent\keywords{Linear Recurrent Neural Networks, State Tracking, DeltaProduct, Householder Products, Adaptive Computation, Group Representations}
